% Zone „Das kennst du schon“ – aus bank/binomische-formeln/zone.jsonl (zusammenbau v0.1)
\einheitenkopf[zone][Kennst du schon]{Das kennst du schon}
\setcounter{aufgabe}{0}
%% TODO Grundvorstellungs-Aufgabe als erste Hauptnummer der Zone (2.8) fehlt in der Bank

%% TODO Titel: Ich-kann-Satz fehlt in der Bank (Platzhalter: Kettenname) – Nr. 1
%% TODO Anweisung: ein Satz über der ersten Teilaufgabe fehlt in der Bank – Nr. 1
\begin{aufgabe}{Multiplizieren mit Vorzeichen und Punkt vor Strich}
%% TODO Darstellung für schwach fehlt in der Bank (binomische-formeln-zone-f1-v1; Abschnitt „Für schwache Schüler“)
\swz{$(-4) \cdot (-8)$ – Ergebnis?}{}
%% TODO Darstellung für schwach fehlt in der Bank (binomische-formeln-zone-f1-v3; Abschnitt „Für schwache Schüler“)
\swz{$(-3) \cdot 5 - 2 \cdot (-7)$ – Ergebnis?}{}
\end{aufgabe}

%% TODO Titel: Ich-kann-Satz fehlt in der Bank (Platzhalter: Kettenname) – Nr. 2
%% TODO Anweisung: ein Satz über der ersten Teilaufgabe fehlt in der Bank – Nr. 2
\begin{aufgabe}{Termwert berechnen, auch mit negativer Zahl, als Probe einer Umformung}
%% TODO Darstellung für schwach fehlt in der Bank (binomische-formeln-zone-f2-v1; Abschnitt „Für schwache Schüler“)
\swz{$3x + 5$ für $x = 4$ – Termwert?}{}
%% TODO Darstellung für schwach fehlt in der Bank (binomische-formeln-zone-f2-v3; Abschnitt „Für schwache Schüler“)
\swz{$2x^2 - 3x$ für $x = -2$ – Termwert?}{}
\end{aufgabe}

%% TODO Titel: Ich-kann-Satz fehlt in der Bank (Platzhalter: Kettenname) – Nr. 3
%% TODO Anweisung: ein Satz über der ersten Teilaufgabe fehlt in der Bank – Nr. 3
\begin{aufgabe}{Gleichartige Glieder mit einer Variablen zusammenfassen, Vorzahl eins beachten, Potenz getrennt halten}
%% TODO Darstellung für schwach fehlt in der Bank (binomische-formeln-zone-f3-v1; Abschnitt „Für schwache Schüler“)
\swz{$4x + 3x$ – zusammengefasst?}{}
%% TODO Darstellung für schwach fehlt in der Bank (binomische-formeln-zone-f3-v3; Abschnitt „Für schwache Schüler“)
\swz{$5x - 8x + 2$ – zusammengefasst?}{}
\end{aufgabe}

%% TODO Titel: Ich-kann-Satz fehlt in der Bank (Platzhalter: Kettenname) – Nr. 4
%% TODO Anweisung: ein Satz über der ersten Teilaufgabe fehlt in der Bank – Nr. 4
\begin{aufgabe}{Zahl mal Klammer und Minusklammer (Distributivgesetz, alle Vorzeichen drehen)}
%% TODO Darstellung für schwach fehlt in der Bank (binomische-formeln-zone-f4-v1; Abschnitt „Für schwache Schüler“)
\swz{$4 \cdot (x + 2)$ – ohne Klammer?}{}
%% TODO Darstellung für schwach fehlt in der Bank (binomische-formeln-zone-f4-v3; Abschnitt „Für schwache Schüler“)
\swz{$-2 \cdot (x - 6)$ – ohne Klammer?}{}
\end{aufgabe}

%% TODO Titel: Ich-kann-Satz fehlt in der Bank (Platzhalter: Kettenname) – Nr. 5
%% TODO Anweisung: ein Satz über der ersten Teilaufgabe fehlt in der Bank – Nr. 5
\begin{aufgabe}{Quadratzahlen bis 15² erkennen und Quadratwurzeln daraus}
%% TODO Darstellung für schwach fehlt in der Bank (binomische-formeln-zone-f5-v1; Abschnitt „Für schwache Schüler“)
\swz{$\sqrt{64}$ – Wert?}{}
%% TODO Darstellung für schwach fehlt in der Bank (binomische-formeln-zone-f5-v3; Abschnitt „Für schwache Schüler“)
\swz{$(-6)^2$ – Wert?}{}
\end{aufgabe}

%% TODO Titel: Ich-kann-Satz fehlt in der Bank (Platzhalter: Kettenname) – Nr. 6
%% TODO Anweisung: ein Satz über der ersten Teilaufgabe fehlt in der Bank – Nr. 6
\begin{aufgabe}{Scheitelpunktform lesen}
%% TODO Darstellung für schwach fehlt in der Bank (binomische-formeln-zone-f6-v1; Abschnitt „Für schwache Schüler“)
\swz{$y = (x - 4)^2 + 1$ – Scheitel?}{}
%% TODO Darstellung für schwach fehlt in der Bank (binomische-formeln-zone-f6-v3; Abschnitt „Für schwache Schüler“)
\swz{$y = (x - 6)^2 - 7$ – Scheitel?}{}
\end{aufgabe}

%% TODO Titel: Ich-kann-Satz fehlt in der Bank (Platzhalter: Kettenname) – Nr. 7
%% TODO Anweisung: ein Satz über der ersten Teilaufgabe fehlt in der Bank – Nr. 7
\begin{aufgabe}{Ausklammern eines gemeinsamen Zahlfaktors oder einer Variablen}
%% TODO Darstellung für schwach fehlt in der Bank (binomische-formeln-zone-f7-v1; Abschnitt „Für schwache Schüler“)
\swz{$5x + 20$ – ausgeklammert?}{}
%% TODO Darstellung für schwach fehlt in der Bank (binomische-formeln-zone-f7-v3; Abschnitt „Für schwache Schüler“)
\swz{$4x^2 - 8x$ – ausgeklammert?}{}
\end{aufgabe}

%% TODO Titel: Ich-kann-Satz fehlt in der Bank (Platzhalter: Kettenname) – Nr. 8
%% TODO Anweisung: ein Satz über der ersten Teilaufgabe fehlt in der Bank – Nr. 8
\begin{aufgabe}{Zahl mal Klammer und Minusklammer (Distributivgesetz, alle Vorzeichen drehen) – Fehler finden}
\swz[3]{Finde den Fehler und rechne richtig: \rechnung{-(x - 9) &= -x - 9}}{}
\end{aufgabe}

%% TODO Titel: Ich-kann-Satz fehlt in der Bank (Platzhalter: Kettenname) – Nr. 9
%% TODO Anweisung: ein Satz über der ersten Teilaufgabe fehlt in der Bank – Nr. 9
\begin{aufgabe}{Zahl mal Klammer und Minusklammer (Distributivgesetz, alle Vorzeichen drehen) – selbst rechnen}
%% TODO Darstellung für schwach fehlt in der Bank (binomische-formeln-zone-f4-v6; Abschnitt „Für schwache Schüler“)
\swz{$-(2x - 5)$ – ohne Klammer?}{}
\end{aufgabe}
